% Place this refinement after the existing three-point reconstruction proof.
\subsection{Two-point entropy identification and generator ambiguity}\label{subsec:two-laplace}

\begin{proposition}[Two arguments for complete triangles]\label{prop:two-laplace}
For every complete $Q\in\C_3$ and any $0<\lambda_1<\lambda_2$, the two full
matrices $\widehat\Psi_Q(\lambda_1),\widehat\Psi_Q(\lambda_2)$ determine all
six rates uniquely within $\C_{\le3}$, without separately calibrating $r,s$.
The inverse is rational and locally Lipschitz at each such source.
\end{proposition}
\begin{proof}
Use $M=(P\widehat\Psi)^{-1}$ and $h=b+d$ from the preceding reconstruction.
The Schur complement gives
\begin{equation}
 M_{xy}(\lambda)=-\frac{ad}{r(\lambda+h)},\qquad
 M_{yx}(\lambda)=-\frac{cb}{s(\lambda+h)}.
 \label{eq:two-offdiagonal}
\end{equation}
Both are nonzero at a complete triangle. Thus
\begin{equation}
 \rho=\frac{M_{xy}(\lambda_1)}{M_{xy}(\lambda_2)}
 =\frac{\lambda_2+h}{\lambda_1+h}>1,
 \qquad h=\frac{\lambda_2-\rho\lambda_1}{\rho-1}.
 \label{eq:two-hidden-escape}
\end{equation}
The two-point formulas for $u_i,v_i$ in the preceding proof now recover
$r,s,a,c$, and Eq.~\eqref{eq:branch} recovers $b,d$. All their denominators
are nonzero at the source. A two-state competitor has zero inverse
off-diagonals; a reciprocal three-state competitor with equal data must
have $ad>0$, which forces its complete support. It is therefore recovered
by the same rational formulas, proving uniqueness in the full stated cap.
\end{proof}

\begin{corollary}[Two arguments identify entropy]\label{cor:two-laplace-entropy}
For every $Q\in\C_{\le3}$ and any $0<\lambda_1<\lambda_2$, the two full
Laplace matrices determine $\sig(Q)$ uniquely within $\C_{\le3}$,
without separately calibrating the observed rates.
\end{corollary}
\begin{proof}
A nonzero off-diagonal of $M=(P\widehat\Psi)^{-1}$ forces complete
three-state support, so Proposition~\ref{prop:two-laplace} recovers $Q$.
Otherwise every compatible model is a reciprocal tree or two-state chain.
Stationarity across each edge cut gives zero current, hence $\sig=0$.
\end{proof}
The zero pattern refers to $M$, equivalently the off-diagonals of
$F=P\widehat\Psi$, not the off-diagonals of $\widehat\Psi$.
Exact entropy identification does not imply continuity at tree boundaries.

Three points remain necessary in the worst case for \emph{generator recovery}
with uncalibrated rates.
At a hidden leaf attached to $x$, $r,s,a,b>0$ and $c=d=0$, so both
off-diagonals in Eq.~\eqref{eq:two-offdiagonal} vanish. The remaining tests are
$G_x(\lambda)=u+v/(\lambda+b)$, with $u=1/r$, $v=a/r>0$, and $G_y=1/s$.
Given two distinct positive arguments, put
\begin{gather}
 D=\frac{G_x(\lambda_1)-G_x(\lambda_2)}{\lambda_2-\lambda_1}>0,\qquad
 v_\beta=D(\lambda_1+\beta)(\lambda_2+\beta),\\
 u_\beta=G_x(\lambda_1)-D(\lambda_2+\beta).
\end{gather}
Every $\beta>0$ with $u_\beta>0$ yields a leaf generator with
\begin{equation}
 (r,s,a,b,c,d)=(u_\beta^{-1},s,v_\beta/u_\beta,\beta,0,0)
 \label{eq:two-leaf-family}
\end{equation}
having the same two full matrices. An open interval around the source
$\beta=b$ is admissible, and $u_\beta'=-D$ makes the generators distinct.
The $y$-leaf is symmetric. Together with Theorem~\ref{thm:inverse}, this
establishes a worst-case minimum of three positive arguments for generator recovery. In the present
reciprocal three-state model with the observed $x$--$y$ pair present, the two
leaf attachments are the only noncomplete irreducible supports; the conclusion
is not a classification of arbitrary sparse larger networks.

For an explicit check, using the tuple order \emph{$(r,s,a,b,c,d)$}, take
\begin{equation}
 \theta_1=(1,1,1,1,0,0),\qquad
 \theta_2=(6/5,1,12/5,2,0,0).
\end{equation}
Direct evaluation of $R(\lambda I-T)^{-1}B$, including the two observed-rate
amplitudes and the $(+,-)$ row order, gives
\begin{equation}
 \widehat\Psi_1(1)=\widehat\Psi_2(1)=
 \begin{pmatrix}0&1/2\\2/5&0\end{pmatrix},\qquad
 \widehat\Psi_1(2)=\widehat\Psi_2(2)=
 \begin{pmatrix}0&1/3\\3/11&0\end{pmatrix}.
 \label{eq:two-leaf-counterexample}
\end{equation}
The complete transforms differ, since
\begin{equation}
 [\widehat\Psi_1-\widehat\Psi_2]_{-,+}(\lambda)
 =-\frac{\lambda(\lambda-1)(\lambda-2)}
 {(\lambda^2+3\lambda+1)(5\lambda^2+28\lambda+12)}.
\end{equation}
At $\lambda=3$ their lower-left entries are $4/19$ and $10/47$.
These are reciprocal trees and both have zero entropy. The example is
ambiguity of the generator from a finite set of summaries, not equality of
the full waiting laws or a nontrivial entropy fiber.

Calibration changes this comparison. The observed rate $r$ is $1$ versus
$6/5$, whereas $s=1$ in both. Their stationary mark intensities in physical
time are $1/3$ per mark versus $6/17$ per mark; their normalized mark
proportions are nevertheless $(1/2,1/2)$ in both. Their traces also differ.
Thus fixed observed rates, a supplied absolute stationary event intensity,
or a supplied trace exclude this particular pair. More generally, if both
$r,s$ are known, subtracting the appropriate constant from $G_i$ lets the
ratio of two nonzero residual tests recover $h$ even at a leaf. Two positive
arguments then suffice for every three-state source in that calibrated class.

The accompanying exact audit retains the one-point generator counterexample
and Jacobian diagnostics. That equilibrium example does not establish that
two arguments are necessary for entropy identification. The three-point
stability theorem retains its separate role near leaf boundaries.
